\section{Related Work}
\label{sec:related}

\para{Benchmarks for false refusal.} A recent line of work measures exaggerated safety by constructing benign prompts that resemble unsafe requests on the surface. \papertitletext{XSTest} introduced a compact suite of harmless prompts with unsafe contrasts \cite{rottger2024xstest}. \phtest expands this idea with pseudo-harmful but benign prompts generated to stress refusal boundaries \cite{an2024phtest}. \orbench scales over-refusal evaluation to a broader prompt distribution \cite{cui2024orbench}, and \papertitletext{COVER} studies the same issue in context-rich settings \cite{sullutrone2025cover}. We build directly on this literature by using four benign benchmark splits to measure held-out refusal after our intervention.

\para{Generalization and mechanisms of refusal.} Several papers suggest that refusal is not tied only to exact prompt text. \papertitletext{Does Refusal Training in LLMs Generalize to the Past Tense?} shows transfer to transformed harmful prompts that preserve semantics while changing surface form \cite{andriushchenko2025pasttense}. \papertitletext{Refusal in Language Models Is Mediated by a Single Direction} argues that refusal behavior is partly organized along a shared latent direction that can be steered or ablated \cite{arditi2024refusal}. These results motivate our hypothesis that even a few benign-refusal examples may generalize beyond their literal wording.

\para{Safety fragility under downstream tuning.} Our study is also related to work showing that alignment can be brittle under narrow post-training updates. \papertitletext{Fine-tuning Aligned Language Models Compromises Safety, Even When Users Do Not Intend To!} finds that benign-looking fine-tuning can still compromise safety behavior \cite{qi2023finetuning}. \papertitletext{Emergent Misalignment} shows that narrow fine-tuning can create broad and unexpected behavioral shifts \cite{betley2025emergent}. We focus on the complementary failure mode: rather than losing harmful refusal, a model can over-generalize refusal and lose benign responsiveness.

\para{Mitigating over-refusal.} Recent mitigation work treats false refusal as a repair problem. Single-vector ablation reduces benign refusal by intervening directly in activation space \cite{wang2024vector}. Reflection-based methods ask models to reason about whether refusal is warranted before answering \cite{si2025think}. The latest trigger-aware analyses frame over-refusal as a consequence of broad refusal cues learned during safety alignment \cite{xue2026deactivating}. Our experiment is upstream of these methods: we study how easily over-refusal can be induced in the first place.

\para{Positioning our work.} Unlike prior benchmarks, we do not only measure false refusal; we deliberately induce it with a tiny supervised fine-tuning intervention. Unlike mechanistic and mitigation studies, we center the causal question of how much benign-refusal supervision is enough to alter a model's held-out behavior. This design lets us test both the strength of generalization and the extent to which the effect is specific to refusal labels rather than generic tiny-data SFT.
